\cleardoublepage
\chapter*{Nota Histórica}
\label{nota-historica}
\addcontentsline{toc}{chapter}{Nota Histórica}
\markboth{Nota Histórica}{}

\thispagestyle{empty}

\noindent Este livro é uma obra de ficção. Joaquim Antônio Correia, o
Paciente 1.260, e a criatura que ele carrega não existiram. A instituição
narrada como \textit{o Colônia} é ficcional.

\vspace{0.8em}

\noindent Mas o solo sobre o qual esta ficção foi construída é real --- e é
por isso que ela precisa terminar aqui, e não com um susto.

\vspace{0.8em}

\noindent Na seca de 1915, o governo do Ceará criou o \textbf{Campo do
Alagadiço}, próximo a Fortaleza, onde retirantes foram cercados e
confinados. O \textbf{Campo do Patu}, em Senador Pompeu, pertence a outro
momento histórico: a seca de 1932, com funcionamento situado em
1932--1933. Esta edição adota Alagadiço como referente do cenário de 1915 e
mantém Joaquim, sua família, a vala e a transferência de 1917 como elementos
ficcionais. Os totais de internados e mortos não são apresentados como
precisos.

\vspace{0.8em}

\noindent O \textbf{Hospital Colônia de Barbacena}, em Minas Gerais,
existiu de verdade. Ao longo do século~XX, estima-se que cerca de
\textbf{60~mil pessoas} morreram entre seus muros --- muitas delas sem
diagnóstico psiquiátrico algum: eram pobres, negras, epilépticas, mães
solteiras, homossexuais, ``inconvenientes''. Houve eletrochoque sem
anestesia, contenção, fome, valas comuns e a venda de corpos a faculdades
de medicina. Nada disso é invenção deste livro; é história documentada ---
ver, entre outras fontes, \textit{Holocausto Brasileiro}, de Daniela Arbex.

\vspace{0.8em}

\noindent O horror sobrenatural destas páginas é uma metáfora. A dor que
atravessa gerações, que ninguém processa e que passa adiante como um
contágio, é a única criatura verdadeira do livro --- e ela tem nome:
esquecimento. O pior que se pode fazer com os mortos dos currais e da
Colônia não é assombrá-los com uma história. É deixá-los sem nenhuma.

\vspace{0.8em}

\noindent \textbf{Sobre o risco de estetização:} os autores têm consciência
de que a literatura do trauma caminha na fronteira entre a elaboração
simbólica e o consumo do sofrimento alheio. Este livro não pretende
oferecer o horror do Colônia como espetáculo, nem substituir o testemunho
das vítimas reais e de seus descendentes --- que permanece insubstituível e
anterior a qualquer ficção. Se a leitura destas páginas produzir apenas
sustos, ela terá falhado; se produzir desconforto diante do esquecimento,
terá cumprido sua função. Ao leitor que se sentir afetado: a interrupção
da leitura é sempre legítima, e o apoio está disponível --- CVV 188
(24 horas) e os CAPS de todo o país.

\vspace{1.2em}

\noindent\hfill\textit{Para que não se apague o que o tempo enterrou.}\hspace{2em}
